\chapter{Modèles linéaires et quadratiques}
\label{manual:chapter:04}

\begin{questionbox}
Comment reconnaître si une situation varie à rythme constant ou si elle possède plutôt un sommet, un maximum ou un minimum?
\end{questionbox}

\section{Le modèle affine}
\label{manual:section:04.01}

\begin{definition}{Fonction affine}{}
Une fonction affine est de la forme
\[
f(x)=mx+b.
\]
Le nombre $m$ est la pente et $b=f(0)$ est l'ordonnée à l'origine.
\end{definition}

La pente entre deux points $(x_1,y_1)$ et $(x_2,y_2)$, avec $x_1\neq x_2$, est
\[
m=\frac{y_2-y_1}{x_2-x_1}.
\]
Elle mesure la variation de $y$ pour une unité de variation de $x$.

\begin{center}
\begin{tikzpicture}[scale=0.8]
\draw[-{Stealth},thick] (-1,0)--(6,0) node[right] {$x$};
\draw[-{Stealth},thick] (0,-1)--(0,5) node[above] {$y$};
\draw[uqamblue,very thick] (-0.3,0.55)--(5.5,4.6);
\fill (1,1.45) circle (2pt) node[below right] {$A$};
\fill (4,3.55) circle (2pt) node[above left] {$B$};
\draw[dashed] (1,1.45)--(4,1.45)--(4,3.55);
\node at (2.5,1.15) {$\Delta x$};
\node at (4.45,2.5) {$\Delta y$};
\end{tikzpicture}
\end{center}

\begin{examplebox}
Un taxi demande $4{,}50\,$\$ au départ et $1{,}80\,$\$ par kilomètre. Le coût en fonction de la distance $d$ est
\[
C(d)=1{,}80d+4{,}50.
\]
La pente $1{,}80$ est le coût marginal par kilomètre; l'ordonnée à l'origine est le coût fixe.
\end{examplebox}

\subsection{Déterminer une droite}

Une droite non verticale est déterminée par deux points, ou par un point et une pente. La forme point-pente est
\[
y-y_0=m(x-x_0).
\]

\begin{examplebox}
La droite passant par $(2,5)$ et de pente $-3$ vérifie
\[
y-5=-3(x-2),
\]
donc $y=-3x+11$.
\end{examplebox}

\section{Le modèle quadratique}
\label{manual:section:04.02}

\begin{definition}{Fonction quadratique}{}
Une fonction quadratique est de la forme
\[
f(x)=ax^2+bx+c,
\qquad a\neq0.
\]
Son graphe est une parabole.
\end{definition}

La forme canonique
\[
f(x)=a(x-h)^2+k
\]
révèle directement le sommet $(h,k)$ et l'axe de symétrie $x=h$.

\begin{center}
\begin{tikzpicture}[scale=0.75]
\draw[-{Stealth},thick] (-4,0)--(5,0) node[right] {$x$};
\draw[-{Stealth},thick] (0,-3)--(0,5) node[above] {$y$};
\draw[teal,very thick,domain=-2.5:4.5,samples=100] plot(\x,{0.45*(\x-1)^2-2});
\draw[dashed,uqamblue] (1,-2.6)--(1,4.5) node[above] {$x=h$};
\fill[deepblue] (1,-2) circle (2.4pt) node[below right] {$(h,k)$};
\end{tikzpicture}
\end{center}

\begin{methodbox}
Pour passer de $ax^2+bx+c$ à la forme canonique :
\begin{enumerate}
\item factoriser $a$ dans les termes contenant $x$;
\item compléter le carré;
\item regrouper les constantes.
\end{enumerate}
Le sommet est aussi obtenu par
\[
h=-\frac{b}{2a},\qquad k=f(h).
\]
\end{methodbox}

\begin{examplebox}
\begin{align*}
f(x)&=2x^2-8x+5\\
&=2(x^2-4x)+5\\
&=2\bigl((x-2)^2-4\bigr)+5\\
&=2(x-2)^2-3.
\end{align*}
Le sommet est $(2,-3)$ et la parabole s'ouvre vers le haut.
\end{examplebox}

\section{Zéros et formule quadratique}
\label{manual:section:04.03}

Les zéros d'une fonction quadratique sont les solutions de
\[
ax^2+bx+c=0.
\]
Le discriminant
\[
\Delta=b^2-4ac
\]
détermine le nombre de solutions réelles : deux si $\Delta>0$, une solution double si $\Delta=0$, aucune si $\Delta<0$.

\begin{theorem}{Formule quadratique}{}
Pour $a\neq0$ et $\Delta=b^2-4ac\geq0$, les solutions réelles de $ax^2+bx+c=0$ sont
\[
x=\frac{-b\pm\sqrt{b^2-4ac}}{2a}.
\]
Lorsque $\Delta<0$, il n'y a pas de solution réelle.
\end{theorem}

\begin{examplebox}
Pour $2x^2-3x-2=0$,
\[
\Delta=(-3)^2-4(2)(-2)=25,
\]
donc
\[
x=\frac{3\pm5}{4}.
\]
Les solutions sont $2$ et $-\frac12$, et
\[
2x^2-3x-2=(2x+1)(x-2).
\]
\end{examplebox}

\section{Modéliser et interpréter}
\label{manual:section:04.04}

\begin{examplebox}
La hauteur d'un objet lancé est modélisée par
\[
h(t)=-4{,}9t^2+19{,}6t+1{,}5.
\]
Le coefficient négatif indique que la parabole s'ouvre vers le bas. Le maximum est atteint à
\[
t=-\frac{19{,}6}{2(-4{,}9)}=2\text{ secondes}.
\]
La hauteur maximale vaut
\[
h(2)=21{,}1\text{ mètres}.
\]
Les racines indiquent les instants où le modèle atteint la hauteur zéro; seule la racine positive possède un sens physique après le lancement.
\end{examplebox}


\section{Voir la variation : différences premières et secondes}
\label{manual:section:04.05}

\begin{visualbox}
Un modèle affine ajoute toujours la même quantité lorsque $x$ augmente d'une unité. Un modèle quadratique n'a pas une variation constante, mais ses \emph{variations} changent à rythme constant. Cette différence apparaît avant même de connaître la formule.
\end{visualbox}

\begin{center}
\begin{tabular}{c@{\hspace{1.4cm}}c}
\begin{tabular}{c|ccccc}
$x$&0&1&2&3&4\\\hline
$2x+1$&1&3&5&7&9\\
$\Delta$&&2&2&2&2
\end{tabular}
&
\begin{tabular}{c|ccccc}
$x$&0&1&2&3&4\\\hline
$x^2$&0&1&4&9&16\\
$\Delta$&&1&3&5&7\\
$\Delta^2$&&&2&2&2
\end{tabular}
\end{tabular}
\end{center}

\begin{center}
\begin{tikzpicture}[scale=0.72]
  \draw[-{Stealth},thick,slate] (-1,0)--(6,0) node[right] {$x$};
  \draw[-{Stealth},thick,slate] (0,-1)--(0,5.2) node[above] {$y$};
  \draw[navy,very thick] (-0.2,0.4)--(5.4,4.6);
  \node[navy] at (4.2,3.2) {variation constante};
\end{tikzpicture}
\qquad
\begin{tikzpicture}[scale=0.72]
  \draw[-{Stealth},thick,slate] (-3,0)--(3.5,0) node[right] {$x$};
  \draw[-{Stealth},thick,slate] (0,-1)--(0,5.2) node[above] {$y$};
  \draw[teal,very thick,domain=-2.6:2.6,samples=80] plot(\x,{0.55*\x*\x+0.3});
  \node[teal] at (1.7,3.7) {sommet et courbure};
\end{tikzpicture}
\end{center}

\begin{connectionbox}
Les trois formes d'un trinôme mettent en évidence trois questions différentes :
\[
ax^2+bx+c \quad\text{(valeur initiale et coefficients)},
\]
\[
a(x-h)^2+k \quad\text{(sommet)},
\qquad
 a(x-r_1)(x-r_2) \quad\text{(zéros)}.
\]
Changer de forme ne change pas la fonction; cela change ce que l'on voit immédiatement.
\end{connectionbox}

\subsection{Le domaine vient du contexte}
Une formule peut être définie pour tout réel alors que le modèle ne l'est pas. Si $t$ représente le temps après un lancement, on impose $t\ge0$; si $h(t)$ représente une hauteur, on ne conserve que la portion où $h(t)\ge0$. Un modèle mathématique est toujours une formule accompagnée de son domaine et de ses unités.


\section{Passer de données à un modèle}
\label{manual:section:04.06}

La première tâche n'est pas d'ajuster une formule compliquée, mais de reconnaître le type de variation présent dans les données.

\subsection{Différences et rapports successifs}
Considérons deux tableaux.
\[
\begin{array}{c|ccccc}
t&0&1&2&3&4\\\hline
A(t)&12&17&22&27&32
\end{array}
\qquad
\begin{array}{c|ccccc}
t&0&1&2&3&4\\\hline
B(t)&12&18&27&40{,}5&60{,}75
\end{array}
\]
Pour $A$, les différences successives sont toutes égales à $5$ : le modèle affine $A(t)=12+5t$ est naturel. Pour $B$, les rapports successifs valent $1{,}5$ : le modèle exponentiel $B(t)=12(1{,}5)^t$ est naturel. Un même jeu de points peut parfois être approximé par plusieurs modèles, mais ces diagnostics indiquent la structure la plus simple.

\subsection{Résidus : regarder ce que le modèle n'explique pas}
Si une observation vaut $y_i$ et que le modèle prédit $\widehat y_i$, le résidu est
\[
r_i=y_i-\widehat y_i.
\]
Un bon modèle ne signifie pas forcément des résidus nuls. Il signifie que les résidus sont petits par rapport à l'échelle du problème et qu'ils ne présentent pas une structure évidente.

\begin{center}
\begin{tikzpicture}[scale=0.82]
 \draw[-{Stealth},thick,slate] (-0.4,0)--(6.5,0) node[right] {$x$};
 \draw[-{Stealth},thick,slate] (0,-0.5)--(0,5.5) node[above] {$y$};
 \draw[navy,very thick] (0.2,0.9)--(6,4.7);
 \foreach \x/\y in {0.8/1.6,1.8/1.75,2.8/3.15,3.8/3.2,4.8/4.55,5.6/4.35}{
   \fill[terracotta] (\x,\y) circle (2.5pt);
   \draw[terracotta,thin,dashed] (\x,\y)--(\x,{0.655*\x+0.77});
 }
 \node[terracotta,align=left] at (5.3,1.3) {écarts verticaux\\= résidus};
\end{tikzpicture}
\end{center}

\subsection{Le sommet comme compromis}
Dans un modèle quadratique $f(x)=a(x-h)^2+k$, le sommet $(h,k)$ donne directement l'optimum. Cette forme est particulièrement naturelle lorsqu'une quantité diminue puis augmente, ou augmente puis diminue : coût total en fonction d'une production, hauteur d'un projectile, aire soumise à une contrainte.

\begin{connectionbox}
Un modèle n'est pas jugé seulement par la proximité numérique. Il doit aussi respecter le contexte : unités correctes, domaine réaliste, comportement plausible et paramètres interprétables. Une parabole peut ajuster quatre points de population sans pour autant décrire une croissance crédible à long terme.
\end{connectionbox}



\subsection{Compléter le carré comme déplacement d'une parabole}
La transformation
\[
x^2+6x+5=(x+3)^2-4
\]
peut être lue comme une opération géométrique. Le graphe de $y=x^2$ est déplacé de $3$ unités vers la gauche et de $4$ unités vers le bas. Le sommet, difficile à lire dans la forme développée, devient immédiatement visible.

\begin{center}
\begin{tikzpicture}[scale=0.72]
 \draw[-{Stealth},thick,slate] (-5,0)--(4.5,0) node[right] {$x$};
 \draw[-{Stealth},thick,slate] (0,-5)--(0,5.5) node[above] {$y$};
 \draw[navy!45,thick,dashed,domain=-2.1:2.1,samples=80] plot(\x,{0.65*\x*\x});
 \draw[teal,very thick,domain=-5:-1,samples=80] plot(\x,{0.65*(\x+3)*(\x+3)-4});
 \fill[terracotta] (-3,-4) circle (3pt);
 \draw[gold,thick,->] (0,0.4)--(-3,0.4);
 \draw[gold,thick,->] (-3,0.2)--(-3,-4);
 \node[gold!75!black,above] at (-1.5,0.4) {$-3$ horizontalement};
 \node[gold!75!black,rotate=90] at (-3.25,-1.8) {$-4$ verticalement};
 \node[teal] at (-1.5,3.4) {$y=(x+3)^2-4$};
\end{tikzpicture}
\end{center}

La même manipulation, appliquée à $ax^2+bx+c=0$, conduit à la formule quadratique. La formule n'est donc pas un résultat extérieur : elle condense une complétion du carré effectuée avec des coefficients généraux.

\subsection{Le discriminant comme distance au niveau zéro}
Dans la forme canonique
\[
f(x)=a(x-h)^2+k,
\]
les racines existent lorsque la parabole atteint le niveau $y=0$. Le discriminant encode cette situation : $\Delta>0$ correspond à deux intersections, $\Delta=0$ à un contact au sommet, et $\Delta<0$ à aucune intersection réelle. Cette lecture graphique précède utilement le calcul des racines.


\section{Construire, comparer et valider un modèle}
\label{manual:section:04.07}

Un modèle mathématique ne reproduit pas toute la réalité. Il retient une structure utile : variation constante pour le modèle affine, courbure régulière et extremum pour le modèle quadratique. La qualité d'un modèle dépend autant de l'interprétation de ses paramètres que de l'exactitude des calculs.

\subsection{Déterminer un modèle affine à partir de deux points}

Une droite passant par $(x_1,y_1)$ et $(x_2,y_2)$, avec $x_1\neq x_2$, a pour pente
\[
a=\frac{y_2-y_1}{x_2-x_1}.
\]
Une fois $a$ connu, l'ordonnée à l'origine est obtenue par $b=y_1-ax_1$.

\begin{examplebox}
Un réservoir contient $620$ L au temps $t=0$ et $470$ L après $3$ h. Si la diminution est approximativement constante,
\[
a=\frac{470-620}{3-0}=-50\ \mathrm{L/h}.
\]
Le modèle est
\[
V(t)=620-50t.
\]
La pente est un taux de variation avec unité; l'ordonnée à l'origine est la quantité initiale.
\end{examplebox}

Le modèle prédit $V(t)=0$ pour $t=12{,}4$ h. Au-delà, il donnerait un volume négatif, ce qui montre que son domaine contextuel doit être limité.

\subsection{Comparer deux modèles affines}

Deux forfaits sont représentés par
\[
A(n)=18+0{,}08n,
\qquad
B(n)=30.
\]
Le seuil d'égalité est déterminé par
\[
18+0{,}08n=30,
\qquad n=150.
\]
Pour $n<150$, le forfait A coûte moins cher; pour $n>150$, le forfait B coûte moins cher. Le point d'intersection du graphique représente donc un seuil de décision.

\subsection{Trois formes d'une fonction quadratique}

Une même fonction quadratique peut être écrite sous trois formes :
\[
f(x)=ax^2+bx+c,
\]
\[
f(x)=a(x-h)^2+k,
\]
\[
f(x)=a(x-r_1)(x-r_2)
\]
lorsque deux racines réelles existent. Chaque forme répond à une question différente :
\begin{itemize}
\item la forme développée montre les coefficients et l'ordonnée à l'origine;
\item la forme canonique montre le sommet et l'axe de symétrie;
\item la forme factorisée montre les zéros et le signe.
\end{itemize}

\begin{examplebox}
Pour
\[
f(x)=2x^2-8x+5,
\]
la complétion du carré donne
\[
f(x)=2(x-2)^2-3.
\]
Le sommet est $(2,-3)$. La forme développée donne immédiatement $f(0)=5$. Les deux représentations décrivent la même parabole, mais rendent visibles des informations différentes.
\end{examplebox}

\subsection{Dérivation de la formule quadratique}

Partons de
\[
ax^2+bx+c=0,
\qquad a\neq0.
\]
En divisant par $a$ et en complétant le carré,
\[
x^2+\frac ba x=-\frac ca,
\]
\[
x^2+\frac ba x+\left(\frac b{2a}\right)^2
=
-\frac ca+\left(\frac b{2a}\right)^2,
\]
\[
\left(x+\frac b{2a}\right)^2
=
\frac{b^2-4ac}{4a^2}.
\]
En multipliant par $4a^2$, on obtient la forme plus directe
\[
(2ax+b)^2=b^2-4ac=\Delta.
\]
Lorsque $\Delta\geq0$,
\[
2ax+b=\pm\sqrt\Delta,
\]
d'où
\[
x=\frac{-b\pm\sqrt{b^2-4ac}}{2a}.
\]
La formule n'est donc pas une règle isolée : elle résume la complétion du carré.

\subsection{Le discriminant comme information graphique}

Le nombre
\[
\Delta=b^2-4ac
\]
indique le nombre d'intersections de la parabole avec l'axe horizontal :
\begin{itemize}
\item $\Delta>0$ : deux racines réelles distinctes;
\item $\Delta=0$ : une racine double, la parabole est tangente à l'axe;
\item $\Delta<0$ : aucune intersection réelle.
\end{itemize}
Le discriminant relie donc l'équation, la factorisation et le graphe.

\subsection{Étude complète d'une trajectoire}

Une balle est lancée depuis $1{,}5$ m avec une vitesse verticale initiale de $14$ m/s :
\[
h(t)=-4{,}9t^2+14t+1{,}5.
\]
Le sommet est atteint lorsque
\[
t_s=-\frac b{2a}=-\frac{14}{2(-4{,}9)}\approx1{,}43\ \mathrm{s}.
\]
La hauteur maximale est
\[
h(1{,}43)\approx11{,}5\ \mathrm{m}.
\]
Le temps d'impact est la racine positive de $h(t)=0$, environ $2{,}96$ s. La racine négative est une solution de l'équation algébrique, mais elle n'appartient pas au domaine temporel du modèle.

Cette distinction entre solution mathématique et solution contextuelle est centrale en modélisation.

\subsection{Optimisation quadratique}

Si $a<0$, la fonction
\[
f(x)=a(x-h)^2+k
\]
possède un maximum $k$ atteint en $x=h$. Si $a>0$, elle possède un minimum. Dans un problème d'aire ou de profit, la forme canonique permet de lire directement l'optimum.

\begin{examplebox}
Un rectangle a un périmètre de $40$ m. Si un côté vaut $x$, l'autre vaut $20-x$ et l'aire est
\[
A(x)=x(20-x)=-x^2+20x.
\]
En complétant le carré,
\[
A(x)=-(x-10)^2+100.
\]
L'aire maximale est $100\ \mathrm{m}^2$, obtenue pour un carré de côté $10$ m.
\end{examplebox}

\subsection{Critères de choix entre modèles affine et quadratique}

Un modèle affine est naturel lorsque les différences successives sont approximativement constantes. Un modèle quadratique est naturel lorsque les premières différences changent régulièrement ou lorsqu'un maximum/minimum est visible. La décision doit aussi considérer le contexte : une croissance réelle peut sembler affine sur un court intervalle sans l'être à long terme.

\begin{center}
\begin{tabularx}{\textwidth}{@{}lYY@{}}
\toprule
Critère & Modèle affine & Modèle quadratique\\\midrule
Variation & additive constante & variation du taux elle-même régulière\\
Graphe & droite & parabole\\
Paramètres visibles & pente, valeur initiale & sommet, ouverture, zéros\\
Usage typique & coût unitaire, vitesse constante & trajectoire, aire, optimisation\\
\bottomrule
\end{tabularx}
\end{center}

\subsection{Contrôler un modèle}

Un modèle doit être vérifié à quatre niveaux :
\begin{enumerate}
\item les unités des paramètres sont cohérentes;
\item les données utilisées sont reproduites;
\item le domaine contextuel est annoncé;
\item les prédictions restent plausibles sur l'intervalle étudié.
\end{enumerate}

\begin{summarybox}
Le modèle affine décrit une variation additive constante; le modèle quadratique décrit une courbure et un extremum. Les différentes formes algébriques rendent visibles des propriétés complémentaires, et le contexte décide quelles prédictions sont admissibles.
\end{summarybox}


